% Paper 1 draft — introduction only. Standalone; preamble matches thesis_outline.tex.
% Status: first draft. Methods and experiments are not in this file.
\documentclass[11pt,a4paper]{article}

\usepackage[margin=1in]{geometry}
\usepackage{setspace}
\usepackage{hyperref}
\usepackage[numbers]{natbib}

\onehalfspacing
\title{Phone-as-Brain Feasibility for a Cheap Differential-Drive Rover\\
\large Introduction (draft)}
\author{Yojan Gautam}
\date{\today}

\begin{document}
\maketitle

\begin{abstract}
\noindent
Draft introduction for the first paper: why a phone-only rover is worth attempting, what feasibility question is still open, and what this paper will contribute.
\end{abstract}

\section{Introduction}
\label{sec:introduction}

Disaster recovery, fire and rescue, and related tasks such as military surveillance spread across ground that is often unsafe for people to occupy for long.
Search, inspection, communications relay, and persistent observation need presence in many places at once.
A fleet of simple vehicles covers more of that ground than one high-end platform, if each vehicle is cheap enough to field, lose, and replace, and simple enough to power, charge, and send out again under time pressure.

Robots that navigate in clutter usually carry a dedicated onboard computer: an industrial PC, a companion single-board computer, or an embedded GPU board, cabled to a separate sensor suite.
That computer is what makes a perception pipeline and a local planner straightforward to host.
It is also a logistics item.
Each unit adds bill-of-materials cost, a thermal and power design, spare parts, a machine image to flash and update, and bring-up steps outside the mechanical assembly.
A lab can absorb that overhead for a single platform.
A first-responder unit that needs many robots pays that overhead on every unit, because each custom computer has to be charged, configured, and diagnosed before the chassis can move.

A current consumer phone already holds the pieces a small rover needs for local navigation: cameras, an inertial measurement unit, radios, a battery, and a CPU and GPU able to run real-time vision, often with a depth-capable sensor as well.
Software reaches the device as an application install.
The premise of this thesis is to use that package as the vehicle's only onboard brain.
An operator downloads an application, mounts the phone on a cheap differential-drive chassis, and the handset senses, computes, and commands the wheels.

Phones already appear in robotics as teleoperation clients, wireless links, and extra cameras.
What this platform still has to establish is closed-loop use on the vehicle itself.
The phone must take images, depth, and inertial measurements, keep a local picture of nearby obstacles, and issue wheel velocities at a rate and latency that let a differential-drive rover follow waypoints and clear clutter.
Planners in the dynamic-window family do this on differential-drive robots: each cycle they score velocities the wheels can execute, trading progress toward a goal against clearance and kinematic limits~\citep{fox1997dwa}.
A phone application can host that class of planner.
Whether the camera pipeline, scheduling, heat, and battery of a consumer handset leave a usable margin on a low-cost chassis, indoors and outdoors, with real depth noise, is settled by trials on the robot.

This paper is the first of three that share one chassis and one application.
It treats a single rover.
Collaborative mapping across phones, and the human interface for supervising a fleet, come later; both assume that one phone-mounted rover can already sense, avoid, and drive.
The research question is the following.

\begin{quote}
Can a consumer phone (application, cameras, IMU, and compute) run the sensing and control loop for a cheap differential-drive rover well enough for waypoint navigation with local obstacle avoidance, without onboard industrial compute?
\end{quote}

The author has built the autonomy stack and the human interface from the ground up.
A rolling chassis with depth sensing and obstacle avoidance is in progress.
This paper states the feasibility claim those experiments are meant to support, and the measurements that decide it.
Kinematics and the planner objective are written in the methods.

The planned contributions are three.

\begin{enumerate}
\item A phone-only sensing and control stack for a differential-drive rover.
  The handset's cameras, IMU, and compute are the sole onboard means of perceiving near-field geometry and producing wheel commands.
\item An experimental feasibility study of waypoint navigation with local obstacle avoidance.
  Reported quantities are control and sensing loop rate, end-to-end latency, waypoint success, obstacle clearance and collisions, human interventions and emergency stops per run, and power draw together with thermal throttling.
\item A cost and setup-time comparison against a dedicated-compute baseline on the same class of chassis, so the result is stated in operational terms as well as in navigation performance.
\end{enumerate}

\bibliographystyle{plainnat}
\begin{thebibliography}{99}

\bibitem[Fox et~al.(1997)]{fox1997dwa}
D.~Fox, W.~Burgard, and S.~Thrun.
\newblock The dynamic window approach to collision avoidance.
\newblock {\em IEEE Robotics \& Automation Magazine}, 4(1):23--33, 1997.

\end{thebibliography}

\end{document}
